纳米硒（SeNPs）细胞毒性预测需要区分同一研究内的拟合表现、跨文献泛化与新材料实验观测，材料化学信息能否在基础描述符之外提供稳定收益也需要独立检验。本研究采用通用纳米材料基准、历史纳米硒分组建模与新增人脐静脉内皮细胞（HUVEC）实验相结合的设计，在 434 条历史记录、24 个文献组上进行 5 折外层、3 折内层嵌套交叉验证，并比较化学增强、剂量—时间交互和固定预算迁移策略。无实测半数抑制浓度（\IC）输入的随机森林基础模型取得折外平均绝对误差 19.11 个百分点、均方根误差 24.12 个百分点和 $R^2=0.109$，新增化学描述符、显式交互项与迁移策略均未显示稳定的误差降低。

独立实验数据包含 9 种 SeNPs 在 HUVEC 中的 351 条报告值，覆盖 24、48、72 h 和 13 个浓度水平。在最高报告浓度 512 下，9 种材料的平均存活率由 24 h 的 59.81\% 降至 72 h 的 26.38\%，材料间差异也随评价条件而变化。27 条材料—时间曲线中，19 条可通过相邻点的对数浓度插值得到表观 \IC，8 条在测试范围内未达到 50\% 存活率。实验提供了剂量—时间响应与材料差异的描述性证据，但现有归档没有事先封存的逐条件模型预测，也未提供原始重复孔数据；浓度单位的含义和计量基准仍待核实。这些结果支持以严格评价和可追溯实验观测约束纳米硒毒性预测的应用范围，尚不足以证明新材料的前瞻预测精度或单一材料性质的因果作用。
